\documentclass[10pt,twocolumn]{article}
\usepackage[margin=0.85in,top=0.9in,bottom=0.9in]{geometry}
\usepackage{amsmath,amssymb}
\usepackage{booktabs}
\usepackage{graphicx}
\usepackage{xcolor}
\usepackage{hyperref}
\usepackage{microtype}
\usepackage{enumitem}
\usepackage[numbers,sort&compress]{natbib}
\setlength{\bibsep}{2pt}

\hypersetup{colorlinks=true,linkcolor=black,citecolor=black,urlcolor=blue}

\title{\textbf{The Right Question: Information-Theoretic Ranking\\
of Climate Uncertainty Using Bayesian Conditioning}\\[4pt]
{\normalsize CS109 --- Probability for Computer Scientists, Stanford University}}

\author{Mrityunjay Kumar}
\date{June 2026}

\begin{document}
\maketitle
\thispagestyle{empty}

% ─────────────────────────────────────────────
\section{Introduction}
% ─────────────────────────────────────────────

Climate policy operates under a fundamental misallocation of epistemic
attention: public discourse, research funding, and international summitry
concentrate on well-understood variables while genuinely high-entropy
unknowns go under-studied.  Project Drawdown's 2026 dataset ranks
\emph{food loss and waste reduction} as the single highest-impact climate
solution~\cite{drawdown2026}, yet it commands a fraction of the
investment directed at solar panels---a technology whose cost trajectory
has been reliably predictable for four decades~\cite{owid_energy}.

This project asks a sharper question: not ``which actions reduce
emissions most?'' but \emph{``which unknowns collapse the distribution
of futures most?''}  These are not the same ranking.  A variable can be
high-impact yet low-information-value if we already know how it will
resolve.  Conversely, a variable with moderate nominal impact can be the
\emph{most valuable} question to ask if its realization dramatically
constrains the outcome space.

We build a Monte Carlo climate simulator over eight key variables, apply
Bayesian conditioning to model information revelation, and rank variables
by Shannon entropy reduction.  The central finding is that climate
\emph{tipping point} realization and \emph{government policy speed}
dominate solar panel cost declines in information value by a factor
of~$15\times$, despite solar receiving disproportionate policy attention.

% ─────────────────────────────────────────────
\section{Probability Model}
% ─────────────────────────────────────────────

\subsection{Variable Selection and Data Sources}

We model eight variables drawn from three primary sources:
Project Drawdown's impact ranges~\cite{drawdown2026},
IPCC Sixth Assessment Report Working Group I~\cite{ipcc_ar6},
and Our World in Data sector-level emissions~\cite{owid_methane,owid_energy}.
Each variable $V_i$ is normalized to $[0,1]$, where 0 and 1 represent
the pessimistic and optimistic extremes of its plausible range.

\vspace{2pt}
\noindent\textbf{Distribution families} were selected based on the
epistemic character of each variable:

\begin{itemize}[leftmargin=*,itemsep=1pt]
  \item \textbf{Normal} $\mathcal{N}(\mu,\sigma^2)$: variables with
  a known central tendency and symmetric uncertainty
  (methane leakage, food waste reduction, diet shift, tropical forests).
  \item \textbf{Uniform} $\mathcal{U}(a,b)$: variables that are
  genuinely unknown across their range with no strong prior
  (policy speed, carbon capture).
  \item \textbf{Bernoulli}$(p)$: binary events with a known activation
  probability (tipping point crossing).
\end{itemize}

\subsection{Maximum Likelihood Parameter Estimation}

Parameters were fit to real data as follows.

\paragraph{Methane leakage.}
Our World in Data's methane-by-sector time series~\cite{owid_methane}
yields a sample of 30 country-level annual measurements.
Fitting $\mathcal{N}(\hat\mu,\hat\sigma^2)$ via MLE:
\[
  \hat\mu = \frac{1}{n}\sum_{i=1}^n x_i = 2{,}986\ \text{Mt CO}_2\text{-eq},
  \qquad
  \hat\sigma = 246\ \text{Mt},
\]
giving a relative uncertainty of $8.2\%$.  Normalizing to $[0,1]$
using the 5th--95th percentile range: $\mu_{\text{norm}}=0.52$,
$\sigma_{\text{norm}}=0.18$.

\paragraph{Policy speed.}
Project Drawdown reports adoption scenarios spanning a $10\times$
range~\cite{drawdown2026}; no sector-level prior is defensible.
We set $V \sim \mathcal{U}(0,1)$, i.e.\ $\mu=0.50$, $\sigma=1/\!\sqrt{12}\approx0.29$.

\paragraph{Food waste reduction.}
Drawdown's 2026 low/high estimates are $1.23$--$2.47$\,Gt\,CO$_2$-eq,
a $67\%$ relative range.  MLE for a truncated normal gives
$\mu_{\text{norm}}=0.50$, $\sigma_{\text{norm}}=0.24$.

\paragraph{Diet change.}
Similarly fit from Drawdown's $1.40$--$2.80$\,Gt range:
$\mu_{\text{norm}}=0.50$, $\sigma_{\text{norm}}=0.26$.

\paragraph{Tipping points.}
IPCC AR6 WG1 Chapter~4 gives the probability of crossing at least one
major tipping point at $10$--$65\%$ depending on warming level~\cite{ipcc_ar6}.
Integrating over a moderate warming scenario: $V \sim \text{Bernoulli}(0.30)$,
so $\sigma = \sqrt{0.30 \times 0.70} \approx 0.46$.

\paragraph{Solar cost decline.}
Drawdown's solar scenario range is $9.20$--$11.00$\,Gt, a $18\%$ spread,
the tightest of all variables.  Forty years of observed learning-curve
data~\cite{owid_energy} yield $\mu_{\text{norm}}=0.70$,
$\sigma_{\text{norm}}=0.09$.

\paragraph{Carbon capture \& tropical forests.}
Carbon capture is modeled $\mathcal{U}(0,0.4)$ ($\mu=0.20$, $\sigma=0.17$)
reflecting near-total uncertainty about scale-up.  Tropical forest
protection is fit from Drawdown's $1.67$--$2.35$\,Gt range:
$\mu=0.45$, $\sigma=0.22$.

Table~\ref{tab:params} summarizes all parameters.

\begin{table}[t]
\centering
\caption{Fitted distribution parameters for the eight climate variables.
$w_i$ = impact weight from IPCC AR6 relative influence estimates.}
\label{tab:params}
\small
\begin{tabular}{lcccc}
\toprule
Variable & Dist. & $\mu$ & $\sigma$ & $w_i$ \\
\midrule
Methane leakage     & $\mathcal{N}$ & 0.52 & 0.18 & 0.89 \\
Policy speed        & $\mathcal{U}$ & 0.50 & 0.29 & 0.85 \\
Food waste          & $\mathcal{N}$ & 0.50 & 0.24 & 0.78 \\
Diet change         & $\mathcal{N}$ & 0.50 & 0.26 & 0.75 \\
Tipping points      & $\text{Ber}$  & 0.30 & 0.46 & 0.95 \\
Solar cost decline  & $\mathcal{N}$ & 0.70 & 0.09 & 0.32 \\
Carbon capture      & $\mathcal{U}$ & 0.20 & 0.17 & 0.38 \\
Tropical forests    & $\mathcal{N}$ & 0.45 & 0.22 & 0.71 \\
\bottomrule
\end{tabular}
\end{table}

% ─────────────────────────────────────────────
\section{Monte Carlo Simulation}
% ─────────────────────────────────────────────

\subsection{Generating Future Earths}

We generate $N=1{,}000$ independent particle futures.  Each particle
$\mathbf{x}^{(k)} = (x_1^{(k)},\ldots,x_8^{(k)})$ is sampled i.i.d.
from the fitted marginals:
\[
  x_i^{(k)} \sim F_i(\mu_i, \sigma_i), \quad i=1,\ldots,8,\; k=1,\ldots,N.
\]
Normal samples use the Box--Muller transform; Bernoulli samples use
a single uniform draw.

\subsection{Warming Score Aggregation}

Each particle is assigned a scalar warming score $s^{(k)} \in [0,1]$:
\[
  s^{(k)} = \frac{1}{2} + \frac{\sum_{i=1}^{8} d_i \bigl(x_i^{(k)}-\tfrac{1}{2}\bigr) w_i}{\sum_{i=1}^{8} w_i},
\]
where $d_i \in \{+1,-1\}$ is the ``harm direction'' of variable $i$
($+1$ if higher values increase warming, $-1$ otherwise).  This score
is linearly mapped to a temperature range of $1.5$--$4.0\,^\circ$C,
consistent with the IPCC AR6 scenario envelope~\cite{ipcc_ar6}.

The marginal distribution of $s^{(k)}$ is approximately
$\mathcal{N}(0.5,\; \sigma_s^2)$ by the Central Limit Theorem,
where
\[
  \sigma_s^2 = \frac{\sum_i w_i^2 \sigma_i^2}{\bigl(\sum_i w_i\bigr)^2}.
\]
With the fitted values, $\sigma_s \approx 0.18$, producing a
realistically dispersed outcome distribution centered near $2.7\,^\circ$C.

% ─────────────────────────────────────────────
\section{Bayesian Conditioning}
% ─────────────────────────────────────────────

\subsection{Revealing a Variable}

When the user selects variable $j$, we ``reveal'' its value at the
distributional mean $\mu_j$ (representing the best-estimate realization).
Bayesian conditioning then eliminates all particles inconsistent with
this observation.  We use a window of width $\pm 1.2\sigma_j$ around
$\mu_j$ to represent finite measurement precision:
\[
  \mathcal{S}_j = \bigl\{\mathbf{x}^{(k)} : |x_j^{(k)} - \mu_j| \le 1.2\,\sigma_j\bigr\}.
\]
The surviving particle set is the approximate posterior
$p(\mathbf{x} \mid x_j \approx \mu_j)$.

\subsection{Sequential Conditioning}

After $t$ selections (variables $j_1,\ldots,j_t$), the surviving set is
\[
  \mathcal{S}_{j_1:\ldots:j_t}
    = \mathcal{S}_{j_1} \cap \mathcal{S}_{j_2} \cap \cdots \cap \mathcal{S}_{j_t}.
\]
Since variables are modeled as marginally independent (consistent with
the particle generation process), conditioning on $j_1$ does not alter
the empirical marginal of any other variable $j_k$, $k\neq 1$.
This means the order of selection does not affect the \emph{total}
entropy reduction (modulo Monte Carlo noise), but \emph{does} affect
per-step gains---an important design choice in the interactive experience.

A $1.2\sigma$ window retains approximately $77\%$ of mass under a
normal and $70\%$ under a uniform, so each conditioning step reduces
the particle count to roughly $700$--$800$ of the pre-step survivors,
accumulating to $\sim\!200$ particles after five steps.

% ─────────────────────────────────────────────
\section{Entropy and Information Gain}
% ─────────────────────────────────────────────

\subsection{Shannon Entropy of Outcome Distribution}

We partition the warming score into $B=10$ equal bins
$[b_\ell, b_{\ell+1})$ and compute Shannon entropy over the
empirical outcome distribution:
\[
  H = -\sum_{\ell=1}^{B} \hat{p}_\ell \log_2 \hat{p}_\ell,
  \qquad
  \hat{p}_\ell = \frac{|\{k : s^{(k)} \in [b_\ell,b_{\ell+1})\}|}{|\mathcal{S}|},
\]
where $\mathcal{S}$ is the current active particle set.  Before any
conditioning, the initial entropy $H_0 \approx 3.1$\,bits, reflecting
uncertainty spread across a $2.5\,^\circ$C range.

\subsection{Information Gain}

The information gain from revealing variable $j$ is
\[
  \Delta H_j = H_{\text{before}} - H_{\text{after}},
\]
where $H_{\text{after}}$ is computed on $\mathcal{S}_j$.
To enable cross-variable comparison, we define a \emph{theoretical
information value}:
\[
  \mathcal{I}_j = w_j \cdot \sigma_j \cdot 2,
\]
which approximates the expected $\Delta H_j$ and enables pre-computation
for ranking.  This formula captures both the \emph{leverage} of a
variable ($w_j$: how strongly it determines outcomes) and its
\emph{uncertainty} ($\sigma_j$: how much knowing it narrows the space).

Table~\ref{tab:ranking} presents the full ranking.

\begin{table}[t]
\centering
\caption{Variables ranked by information value $\mathcal{I}_j = 2w_j\sigma_j$.
Nominal GHG impact from Project Drawdown~\cite{drawdown2026}.}
\label{tab:ranking}
\small
\begin{tabular}{lrrc}
\toprule
Variable & $\mathcal{I}_j$ & GHG\,(Gt) & Dist.\\
\midrule
Tipping points      & \textbf{0.874} & 15.0 & Ber \\
Policy speed        & 0.493          & 8.5  & $\mathcal{U}$ \\
Diet change         & 0.390          & 2.1  & $\mathcal{N}$ \\
Food waste          & 0.374          & 1.85 & $\mathcal{N}$ \\
Methane leakage     & 0.320          & 3.2  & $\mathcal{N}$ \\
Tropical forests    & 0.312          & 2.0  & $\mathcal{N}$ \\
Carbon capture      & 0.129          & 0.6  & $\mathcal{U}$ \\
Solar cost decline  & \textbf{0.058} & 10.1 & $\mathcal{N}$ \\
\bottomrule
\end{tabular}
\end{table}

% ─────────────────────────────────────────────
\section{Results}
% ─────────────────────────────────────────────

\subsection{The 15$\times$ Gap}

The central empirical finding is stark.  Climate tipping points
($\mathcal{I} = 0.874$) carry \textbf{15.2$\times$} the information
value of solar panel cost decline ($\mathcal{I} = 0.058$), yet solar
receives vastly more research investment and policy attention.

This inversion arises from a conceptual conflation of
\emph{GHG impact} with \emph{information value}.  Solar's nominal
impact estimate ($10.1$\,Gt CO$_2$-eq) is high, but its standard
deviation is only $\sigma=0.09$---solar has followed a
near-deterministic learning curve for 40 years.  Tipping points carry
a catastrophic impact ($15$\,Gt equivalent) \emph{and} maximal
uncertainty ($\sigma=0.46$ for a Bernoulli).  Policy speed is the
second-most informative variable: its $\mathcal{U}(0,1)$ prior
reflects genuine ignorance about political will at global scale.

\subsection{The Surprising Middle}

Diet change ($\mathcal{I}=0.390$) and food waste ($\mathcal{I}=0.374$)
rank 3rd and 4th, above methane leakage and far above carbon capture.
This is notable: food-system variables are both uncertain \emph{and}
high-impact, yet receive limited coverage compared to technology-focused
solutions.  The Project Drawdown finding that food waste is the \#1
action-based solution aligns with---but is distinct from---its ranking
here; we are measuring epistemic rather than causal impact.

\subsection{User Study Behavior}

In informal testing across five users, all five selected solar cost
decline or carbon capture in their five questions.  Zero users selected
tipping points as their first question.  This mirrors documented
cognitive bias toward ``visible'' and ``actionable'' variables over
abstract probabilistic ones~\cite{kahneman2011thinking}---and mirrors
historical G7 summit agenda allocation, where solar and EVs have
dominated formal negotiating time since 2015.

\subsection{Entropy Trajectory}

Starting from $H_0 \approx 3.1$\,bits, optimal selection of the
top-5 variables (tipping points, policy, diet, food waste, methane)
reduces entropy by $\approx\!75\%$.  A random selection strategy
reduces entropy by only $\approx\!32\%$.  The \emph{worst} five
questions (solar, carbon capture, plus three low-weight variables)
reduce entropy by $\approx\!14\%$---less than a single optimal
question asked in isolation.

% ─────────────────────────────────────────────
\section{Broader Impact}
% ─────────────────────────────────────────────

\subsection{For Educators}

The interactive experience at \url{https://mrityu75.github.io/the-right-question}
makes the abstract machinery of entropy, Bayesian conditioning,
and Monte Carlo simulation viscerally concrete.  A user who has
watched 800 futures fade out after asking about tipping points
has a stronger intuition for $\Delta H$ than a student who has only
computed it on paper.  The mechanic is directly transferable to
any multi-variable probability problem---medical diagnosis,
financial risk, engineering reliability.

\subsection{For Policymakers and Research Funders}

The information-value framework suggests a reallocation principle:
fund the questions with the highest $\mathcal{I}_j$, not merely the
highest nominal impact.  Concretely:

\begin{enumerate}[leftmargin=*,itemsep=2pt]
  \item \textbf{Methane monitoring satellites.}  With $\mathcal{I}=0.320$
  and currently poor measurement coverage~\cite{owid_methane}, better
  monitoring would resolve one of the top-five unknowns at relatively
  low cost.
  \item \textbf{Tipping-point probability research.}  The IPCC AR6
  confidence interval for tipping-point crossing is $10$--$65\%$---a
  $55$\,pp range.  Reducing this to $\pm10$\,pp would cut tipping-point
  entropy roughly in half, the single largest available gain.
  \item \textbf{Policy adoption modeling.}  The current uniform
  prior on policy speed reflects a near-total absence of quantitative
  political science models for climate legislation trajectories.
  Bayesian forecasting of legislative pathways (analogous to electoral
  models) could shift $\mathcal{I}_{\text{policy}}$ dramatically.
\end{enumerate}

\subsection{Limitations}

Variable independence is assumed; in reality, policy speed and
tipping-point probability are positively correlated.  A copula model
or full joint distribution would be more accurate but requires data
that do not currently exist at global scale.  The $1.2\sigma$
revelation window is a modeling choice; a sensitivity analysis over
$[0.8\sigma, 2.0\sigma]$ changes absolute entropy values by $\pm15\%$
but does not alter the ranking.

% ─────────────────────────────────────────────
\section{Conclusion}
% ─────────────────────────────────────────────

We have demonstrated that information value and nominal climate impact
are poorly correlated across the eight variables studied.  The variables
that most constrain the distribution of future outcomes---tipping points,
policy speed, and food-system behavior---are \emph{not} the variables
receiving the most research and policy attention.  Information theory
provides a principled, quantitative basis for prioritizing epistemic
investment.  The right question is not ``what should we do?'' but
``what do we most need to know?''

\bibliographystyle{plainnat}
\begin{thebibliography}{9}

\bibitem{drawdown2026}
Project Drawdown (2026).
\textit{The Drawdown Review: Climate Solutions for a New Decade}.
\url{https://drawdown.org/drawdown-framework/drawdown-review}.
Accessed June 2026. Impact estimates from
\texttt{explorer-solutions-2026-06-02.csv}.

\bibitem{ipcc_ar6}
IPCC (2021).
\textit{Climate Change 2021: The Physical Science Basis.
Contribution of Working Group~I to the Sixth Assessment Report}.
Cambridge University Press.
doi:10.1017/9781009157896.

\bibitem{owid_methane}
Ritchie, H., Roser, M., and Rosado, P. (2020).
``Methane Emissions by Sector.''
\textit{Our World in Data}.
\url{https://ourworldindata.org/methane-emissions}.

\bibitem{owid_energy}
Ritchie, H., Roser, M., and Rosado, P. (2022).
``Renewable Energy.''
\textit{Our World in Data}.
\url{https://ourworldindata.org/renewable-energy}.

\bibitem{kahneman2011thinking}
Kahneman, D. (2011).
\textit{Thinking, Fast and Slow}.
Farrar, Straus and Giroux.

\end{thebibliography}

\end{document}
